
% SGA 3，第 VII 篇_B；英文重建本第一轮的中文译本。
% 范围： 节 3.3.2 全文 依据文献局部页 72;
% 止于下节之前： 节 3.3.3.
% 合并阅读本页： 585. 重版印刷页： 574.
% 范围内原版边码： 539.
% 范围内编者注： 131.
% 依据： Polo--Gille Exp7B-13oct24.pdf,
% SHA-256 D303690E0195654BBE7C3DAA5D89EA23E24CD709A3575217234AC19AE8836694.
% 参照了 Jacob Reinhold 的 CC BY 4.0 Markdown，作为对照来源。
% 本组件无图；所有数学内容均为原生 LaTeX。
\subsubsection*{3.3.2 推论}
\phantomsection
\label{sga3:VIIB:3.3.2}
\noindent\textbf{推论 3.3.2。}
{\itshape 设 \(k\) 是特征 \(0\) 的域。函子
\[
  \mathbf L\longmapsto\mathscr G(\mathbf L)
\]
给出 \(k\)-Lie 代数范畴与无穷小形式 \(k\)-群范畴之间的等价。}
\footnote{\label{sga3:VIIB:nde:131}由 \SGARefLink{sga3:VIIB:2.7:prop}{2.7}、\SGARefLink{sga3:VIIB:2.2.1}{2.2.1} 和 \SGARefLink{sga3:VIIB:1.3.6}{1.3.6}，无穷小形式 \(k\)-群 \(G\) 与连通余交换 \(k\)-双代数 \(H\)“是一回事”\textup{（参见\ \SGARefLink{sga3:VIIB:2.9}{2.9}）}。因此，这个推论等价于下面的定理，Kostant 也曾独立得到该定理\textup{（参见\ \SGARefLink{sga3:VIIB:2.9.2}{2.9.2}）}。J.~W.~Milnor 和 J.~C.~Moore 更早在\textup{[MM65]} 中，在附加假设 \(H\) 作为代数由其本原元素生成的条件下，证明了该定理。虽然发表于 1965 年，他们的文本在 1960 年以前就已流传\textup{（参见评论 [Ja65]）}；因此这个结果常称为 Cartier--Kostant--Milnor--Moore 定理。
\smallskip
\noindent\textbf{Cartier--Kostant--Milnor--Moore 定理。}
设 \(k\) 是特征 \(0\) 的域。函子
\[
  \mathbf L\longmapsto U(\mathbf L),
  \qquad
  H\longmapsto\operatorname{Prim}(H),
\]
定义 \(k\)-Lie 代数范畴与连通余交换 \(k\)-双代数范畴之间的等价。
\smallskip
若 \(k\) 是包含 \(\mathbb Q\) 的 Artin 环，则 \SGARefLink{sga3:VIIB:3.2:prop}{3.2}\textup{\SGARefLink{sga3:vii:paragraph:5-1:equation:ii}{(ii)}} 和 \SGARefLink{sga3:VIIB:3.3:cartier-theorem}{3.3}，结合 \SGARefLink{sga3:VIIB:2.7:prop}{2.7}、\SGARefLink{sga3:VIIB:2.2.1}{2.2.1} 和 \SGARefLink{sga3:VIIB:1.3.6}{1.3.6}，同样表明
\[
  \mathbf L\longmapsto\mathscr G(\mathbf L)
\]
\textup{（分别为 \(\mathbf L\mapsto U(\mathbf L)\)）}是从在 \(k\) 上平坦的 \(k\)-Lie 代数范畴到拓扑平坦无穷小形式 \(k\)-群范畴\textup{（分别到在 \(k\) 上平坦的连通余交换 \(k\)-双代数范畴）}的等价。}
事实上，当 \(G\) 遍历无穷小形式 \(k\)-群时，\SGARefLink{sga3:VIIB:3.3:cartier-theorem}{定理 3.3} 将函子
\[
  G\longmapsto\mathscr G\bigl(\operatorname{Lie}G\bigr)
\]
与恒等函子等同。同样，\SGARefLink{sga3:VIIB:3.2:prop}{3.2}\textup{\SGARefLink{sga3:vii:paragraph:5-1:equation:ii}{(ii)}} 将
\[
  \mathbf L
  \longmapsto
  \operatorname{Lie}\bigl(\mathscr G(\mathbf L)\bigr)
  =
  \operatorname{Prim}U(\mathbf L)
\]
与 \(k\)-Lie 代数上的恒等函子等同。
